\chapter{Use Cases and Indexing Solutions}
\label{ch:solutions}
% Small access-path sketches. Boxes represent stored records or operations, not timing claims.
\tikzset{record/.style={draw=blue!65!black,fill=blue!5,rounded corners=2pt,inner sep=5pt,font=\footnotesize,align=center},
  chosen/.style={record,fill=green!12,draw=green!50!black},
  excluded/.style={record,fill=red!8,draw=red!55!black},
  flow/.style={-{Stealth[length=2mm]},thick,blue!65!black}}

\newcommand{\UCOneDiagram}{%
\begin{figure}[H]\centering
\begin{tikzpicture}[x=1cm,y=1cm]
\node[record] at (1.6,1.55) {\texttt{CTXSYS | INDEX | 19}};
\node[chosen] at (1.6,0.85) {\texttt{CTXSYS | TABLE | 12}};
\node[record] at (1.6,0.15) {\texttt{SYS | INDEX | 8}};
\node[chosen] at (1.6,-0.55) {\texttt{SYS | TABLE | 16}};
\node[record] (prefix) at (7,1.4) {\textbf{Leading prefix}\\\texttt{OWNER=CTXSYS}\\read TYPE, ID in order};
\node[chosen] (full) at (7,-0.15) {\textbf{Full key}\\\texttt{SYS, TABLE, 16}\\one narrow range};
\node[excluded] at (7,-1.55) {\textbf{Missing OWNER}\\TABLE entries occur in many owner ranges};
\draw[flow] (3.1,0.85) -- (prefix.west); \draw[flow] (3.1,-0.55) -- (full.west);
\end{tikzpicture}
\caption{UC1: lexicographically ordered compound keys. A leading owner narrows the range; trailing keys already supply the sort order. Missing the owner prevents one contiguous seek.}\label{fig:diag-uc1}
\end{figure}}

\newcommand{\UCTwoDiagram}{%
\begin{figure}[H]\centering
\begin{tikzpicture}[x=1cm,y=1cm]
\node[record,minimum width=3.4cm] (base) at (0,0) {Base rows\\\texttt{N  N  Y  N  Y  N}};
\node[record,minimum width=3.4cm] (full) at (5,0.95) {Full index B\\\texttt{N  N  N  N  Y  Y}};
\node[chosen,minimum width=3.4cm] (part) at (5,-0.95) {Filtered index C\\\texttt{Y  Y}};
\draw[flow] (base.north east) -- ++(0.55,0) |- (full.west);
\draw[flow] (base.south east) -- ++(0.55,0) |- (part.west);
\end{tikzpicture}
\caption{UC2: B stores every row; C stores only Y rows. A query for N cannot use C as a complete access path.}\label{fig:diag-uc2}
\end{figure}}

\newcommand{\UCThreeDiagram}{%
\begin{figure}[H]\centering
\begin{tikzpicture}[node distance=5mm and 8mm]
\node[record] (n) {Narrow B\\\texttt{OBJECT\_ID + pointer}};
\node[excluded,right=of n] (lookup) {Base-table lookup\\once per matching row};
\node[record,right=of lookup] (result) {Projected result};
\node[chosen,below=9mm of n] (c) {Covering C\\\texttt{OBJECT\_ID + output columns}};
\draw[flow] (n) -- (lookup); \draw[flow] (lookup) -- (result);
\draw[flow] (c.east) -- ++(6.3,0) |- (result.south);
\draw[red!60!black,thick] (lookup.south west) -- (lookup.north east);
\end{tikzpicture}
\caption{UC3: C stores the projected columns and bypasses repeated base-row lookup. Wider entries increase index size.}\label{fig:diag-uc3}
\end{figure}}

\newcommand{\UCFourDiagram}{%
\begin{figure}[H]\centering
\begin{tikzpicture}[node distance=5mm and 8mm]
\node[record] (doc) {Article BODY\\\emph{climate change is urgent}};
\node[record,right=of doc] (tok) {Tokenization\\\texttt{climate | change | urgent}};
\node[chosen,right=of tok] (post) {Posting lists\\\texttt{climate} $\to$ IDs 12, 89\\\texttt{change} $\to$ IDs 12, 91};
\node[record,right=of post] (ids) {Candidate IDs\\phrase check if needed};
\draw[flow] (doc) -- (tok); \draw[flow] (tok) -- (post); \draw[flow] (post) -- (ids);
\end{tikzpicture}
\caption{UC4: an inverted index maps words to document IDs. Substring scans inspect characters; linguistic search uses token boundaries, so their ID sets can differ.}\label{fig:diag-uc4}
\end{figure}}

\newcommand{\UCFiveDiagram}{%
\begin{figure}[H]\centering
\begin{tikzpicture}[x=1cm,y=1cm]
\node[record,minimum width=3.5cm] (row) at (0,0.9) {Rowstore B\\\texttt{[TYPE, ID, OWNER, ...]}\\\texttt{[TYPE, ID, OWNER, ...]}};
\node[chosen,minimum width=3.5cm] (col) at (0,-0.9) {Columnstore C\\\textbf{TYPE:} \texttt{TABLE, VIEW, ...}\\\textbf{ID:} \texttt{16, 17, ...}\\OWNER: \texttt{SYS, SYS, ...}};
\node[chosen,minimum width=2.6cm] (agg) at (6,-0.9) {\texttt{GROUP BY TYPE}\\\texttt{SUM(ID)}};
\draw[flow] (col.east) -- node[above=2mm,font=\scriptsize]{two segments} (agg.west);
\end{tikzpicture}
\caption{UC5: C groups and compresses values by column. This aggregation reads TYPE and ID segments; a point predicate without a B-tree still scans.}\label{fig:diag-uc5}
\end{figure}}

\newcommand{\UCSixDiagram}{%
\begin{figure}[H]\centering
\begin{tikzpicture}[node distance=6mm and 10mm]
\node[record] (doc) {Document 127\\\texttt{tags=[owner:SYS,}\\\texttt{type:TABLE, temporary:Y]}};
\node[chosen,right=of doc] (entries) {Multikey entries\\\texttt{owner:SYS $\to$ 127}\\\texttt{type:TABLE $\to$ 127}\\\texttt{temporary:Y $\to$ 127}};
\node[record,right=of entries] (fetch) {Two-tag query\\seek selective tag\\FETCH document\\check second tag};
\draw[flow] (doc) -- (entries); \draw[flow] (entries) -- (fetch);
\end{tikzpicture}
\caption{UC6: one array yields several index entries. A two-tag query may seek one entry then fetch the document to verify the second tag.}\label{fig:diag-uc6}
\end{figure}}


This chapter explains the six database problems we want to solve, the queries that represent them, why a specific indexing mechanism is selected for each workload, and what theoretical trade-off that index introduces. No benchmark numbers are presented yet; our goal here is to establish the engineering logic before running experiments.

\section{UC1: Multi-Column Lookups and Sorting \texorpdfstring{$\rightarrow$}{->} Compound B-Tree Index}

\ReportBlock{Problem}
In large metadata catalogs, users rarely search by a single column. Instead, they filter objects by combinations of \texttt{OWNER}, \texttt{OBJECT\_TYPE}, and \texttt{OBJECT\_ID}, or ask for all objects of an owner sorted by \texttt{(OBJECT\_TYPE, OBJECT\_ID)}. If we only build a single-column index on \texttt{OWNER}, any query for a common owner (such as \texttt{OWNER = 'SYS'}) still scans tens of thousands of index entries and must sort the results afterward.

\Needspace{15\baselineskip}
\ReportBlock{Representative Query}
The full-key lookup (U1a) projects the same four fields in both systems. The other variants are U1b (two-column prefix), U1c (missing \texttt{OWNER}), and U1d (\texttt{OWNER=CTXSYS} sorted by \texttt{OBJECT\_TYPE, OBJECT\_ID}).
\ReportBlock{SQL Server --- U1a}
\begin{SQLCode}
SELECT ID, OWNER, OBJECT_TYPE, OBJECT_ID
FROM dbo.UC1_Bench_100K
WHERE OWNER = N'SYS' AND OBJECT_TYPE = N'TABLE'
  AND OBJECT_ID = 16;
\end{SQLCode}
\ReportBlock{MongoDB --- U1a}
\begin{MongoCode}
db.uc1_bench_100K.find(
  {OWNER: 'SYS', OBJECT_TYPE: 'TABLE', OBJECT_ID: 16},
  {_id: 0, ID: 1, OWNER: 1, OBJECT_TYPE: 1, OBJECT_ID: 1}
);
\end{MongoCode}

\ReportBlock{Chosen Index}
We compare a conventional single-column index on \texttt{(OWNER)} (\textbf{Config B}) against a \textbf{Compound B-Tree Index} on \texttt{(OWNER, OBJECT\_TYPE, OBJECT\_ID)} (\textbf{Config C}) \cite{mssqlindex,mongocompound}. Figure~\ref{fig:diag-uc1} shows its lexicographic order: entries are grouped first by \texttt{OWNER}, then by \texttt{OBJECT\_TYPE}, then by \texttt{OBJECT\_ID}.

\ReportBlock{Representative index definitions --- Config C}
\begin{SQLCode}
CREATE INDEX IX_UC1_OwnerTypeObjectId
ON dbo.UC1_Bench_100K (OWNER, OBJECT_TYPE, OBJECT_ID);
\end{SQLCode}
\begin{MongoCode}
db.uc1_bench_100K.createIndex(
  {OWNER: 1, OBJECT_TYPE: 1, OBJECT_ID: 1},
  {name: 'ix_uc1_owner_type_object_id'}
);
\end{MongoCode}

\ReportBlock{Why It Works}
When a query provides the leading columns (\texttt{U1a}, \texttt{U1b}), the DBMS can seek directly to that narrow range in the tree. When a query filters by the first column (\texttt{OWNER = 'CTXSYS'}) and sorts by the remaining columns (\texttt{U1d}), the leaf pages are already stored in the exact \texttt{ORDER BY} sequence, eliminating the blocking \texttt{Sort} operator. However, if a query omits the first column \texttt{OWNER} (\texttt{U1c}), the matching \texttt{OBJECT\_TYPE = 'TABLE'} entries are scattered across different owners, so the DBMS cannot perform a single range seek. \ReportBlock{Trade-Off}
The main trade-off is that storing three columns in every index entry makes the index larger than a single-column index and adds overhead whenever any of the three columns is modified.

\UCOneDiagram

\section{UC2: Skewed Active-Flag Queries \texorpdfstring{$\rightarrow$}{->} Filtered and Partial Indexes}

\ReportBlock{Problem}
Operational tables often include a status or flag column where almost all rows are inactive (\texttt{TEMPORARY = 'N'}, $99.76\%$) and only a tiny fraction are active (\texttt{TEMPORARY = 'Y'}, $0.24\%$). Application dashboards frequently query the active \texttt{'Y'} subset, while batch jobs occasionally query \texttt{'N'}. Building a full index across the entire table wastes over $99\%$ of index space on \texttt{'N'} rows that active-subset queries never touch.

\Needspace{15\baselineskip}
\ReportBlock{Representative Query}
Both systems return \texttt{OBJECT\_TYPE, OBJECT\_ID, TEMPORARY} for the rare Y subset. The N variant changes only the flag; a separate SQL probe passes it as a parameter.
\ReportBlock{SQL Server --- Y subset}
\begin{SQLCode}
SELECT OBJECT_TYPE, OBJECT_ID, TEMPORARY
FROM dbo.UC2_Bench_1000K
WHERE TEMPORARY = N'Y' AND OBJECT_TYPE = N'TABLE'
  AND OBJECT_ID BETWEEN 1 AND 10000;
\end{SQLCode}
\ReportBlock{MongoDB --- Y subset}
\begin{MongoCode}
db.uc2_bench_1000K.find(
  {TEMPORARY: 'Y', OBJECT_TYPE: 'TABLE',
   OBJECT_ID: {$gte: 1, $lte: 10000}},
  {_id: 0, OBJECT_TYPE: 1, OBJECT_ID: 1, TEMPORARY: 1}
);
\end{MongoCode}

\ReportBlock{Chosen Index}
We compare a full composite index on all rows (\textbf{Config B}) against a \textbf{Filtered Index} in SQL Server (\texttt{WHERE TEMPORARY = N'Y'}) and a \textbf{Partial Index} in MongoDB (\texttt{partialFilterExpression: \{TEMPORARY: "Y"\}}) (\textbf{Config C}) \cite{mssqlfiltered,mongopartial}. As illustrated in Figure~\ref{fig:diag-uc2}, the DBMS inserts a row into Index C only if it satisfies \texttt{TEMPORARY = 'Y'}.

\ReportBlock{Representative index definitions --- Config C}
\begin{SQLCode}
CREATE INDEX IX_UC2_Filtered ON dbo.UC2_Bench_1000K
  (OBJECT_TYPE, OBJECT_ID) INCLUDE (TEMPORARY)
WHERE TEMPORARY = N'Y';
\end{SQLCode}
\begin{MongoCode}
db.uc2_bench_1000K.createIndex(
  {OBJECT_TYPE: 1, OBJECT_ID: 1, TEMPORARY: 1},
  {name: 'ix_uc2_partial',
   partialFilterExpression: {TEMPORARY: 'Y'}}
);
\end{MongoCode}

\ReportBlock{Why It Works}
Because only $0.24\%$ of rows qualify, Index C is much smaller than a full index and does not add a new entry for a row that remains \texttt{'N'}. \ReportBlock{Trade-Off}
The trade-off is strict query eligibility: if a query asks for \texttt{TEMPORARY = 'N'}, Index C contains none of the needed rows and the DBMS must fall back to a full table scan. A parameterized SQL predicate may not prove Y-only eligibility to the optimizer. Both cached and recompiled parameter probes scanned in this local run; the actual plan must be checked.



\UCTwoDiagram

\section{UC3: Range Projections \texorpdfstring{$\rightarrow$}{->} Covering Index (\texttt{INCLUDE} vs.\ Compound Keys)}

\ReportBlock{Problem}
Applications rarely select only the column in the \texttt{WHERE} clause; they usually project several descriptive columns for display. When a query filters by a range on \texttt{OBJECT\_ID} and returns four other columns, a narrow index on \texttt{(OBJECT\_ID)} can find the matching row pointers, but it must jump back to the base table (\texttt{RID Lookup} in SQL Server or \texttt{FETCH} in MongoDB) once for every matched row. When hundreds or thousands of rows match, these random base-table lookups become slower than scanning the entire table.

\Needspace{15\baselineskip}
\ReportBlock{Representative Query}
At 1M rows, the medium range contains 2,775 results; the 100K medium bound is 1,000. The point query uses \texttt{OBJECT\_ID=16} with \texttt{ID=1}; small ranges end at 50 (100K) or 5 (1M).
\ReportBlock{SQL Server --- 1M medium range}
\begin{SQLCode}
SELECT ID, OWNER, OBJECT_NAME, OBJECT_TYPE
FROM dbo.UC3_Bench_1000K
WHERE OBJECT_ID BETWEEN 1 AND 200;
\end{SQLCode}
\ReportBlock{MongoDB --- 1M medium range}
\begin{MongoCode}
db.uc3_bench_1000K.find(
  {OBJECT_ID: {$gte: 1, $lte: 200}},
  {_id: 0, ID: 1, OWNER: 1,
   OBJECT_NAME: 1, OBJECT_TYPE: 1}
);
\end{MongoCode}

\ReportBlock{Chosen Index}
We compare a narrow index on \texttt{(OBJECT\_ID)} (\textbf{Config B}) against a \textbf{Covering Index} (\textbf{Config C}) that contains all columns required by the query \cite{mssqlinclude,mongoexplain}. Figure~\ref{fig:diag-uc3} shows the saved base-row lookup. SQL Server stores the four included columns at the leaf level, keeping upper branch nodes narrow. MongoDB appends all returned fields as compound keys and excludes \texttt{\_id} from the projection. The native definitions below show the distinct mechanisms.

\ReportBlock{Representative index definitions --- Config C}
\begin{SQLCode}
CREATE INDEX IX_UC3_Cover ON dbo.UC3_Bench_1000K
  (OBJECT_ID) INCLUDE (ID, OWNER, OBJECT_NAME, OBJECT_TYPE);
\end{SQLCode}
\begin{MongoCode}
db.uc3_bench_1000K.createIndex(
  {OBJECT_ID: 1, ID: 1, OWNER: 1,
   OBJECT_NAME: 1, OBJECT_TYPE: 1},
  {name: 'ix_uc3_cover'}
);
\end{MongoCode}

\ReportBlock{Why It Works}
Because the index leaves already hold every requested column, the DBMS answers the query directly from the index without touching the base table (\texttt{RID Lookup} disappears in SQL Server; \texttt{totalDocsExamined = 0} in MongoDB). \ReportBlock{Trade-Off}
The trade-off is a substantial increase in index size, because four wide columns are duplicated inside the index, and any update to those extra columns must now also update the index.

\UCThreeDiagram

\section{UC4: Natural-Language Search over Prose \texorpdfstring{$\rightarrow$}{->} Full-Text and Text Indexes}

\ReportBlock{Problem}
Standard B-trees sort strings from their first character. They cannot efficiently find a word or phrase buried in the middle of a multi-paragraph news article. Without a specialized index, searching article text (\texttt{BODY}) requires scanning every character of every document using \texttt{LIKE '\%term\%'} or \texttt{\$regex}.

\Needspace{15\baselineskip}
\ReportBlock{Representative Query}
The six searches comprise three words (election, climate, bitcoin) and three phrases (white house, climate change, stock market). The phrase below shows substring A and linguistic C over the same \texttt{BODY} field; their returned ID sets need not match.
\ReportBlock{SQL Server --- phrase baseline and full-text}
\begin{SQLCode}
SELECT DOC_ID FROM dbo.UC4_Read_A
WHERE BODY LIKE N'%climate change%';
SELECT DOC_ID FROM dbo.UC4_Read_C
WHERE CONTAINS(BODY, '"climate change"', LANGUAGE 1033);
\end{SQLCode}
\ReportBlock{MongoDB --- phrase baseline and text search}
\begin{MongoCode}
db.uc4_read_A.find(
  {BODY: {$regex: 'climate change', $options: 'i'}},
  {_id: 0, DOC_ID: 1}
);
db.uc4_read_C.find(
  {$text: {$search: '"climate change"'}},
  {_id: 0, DOC_ID: 1}
);
\end{MongoCode}

\ReportBlock{Chosen Index}
We build a \textbf{Full-Text Index} in SQL Server and a \textbf{Text Index} in MongoDB (\textbf{Config C}) \cite{mssqlfts,mongotext}. Unlike a B-tree over the whole string, these indexes map tokens to posting lists of document IDs (Figure~\ref{fig:diag-uc4}). Language processing may remove common words or stem related forms; exact behavior depends on the DBMS and query.

\ReportBlock{Representative index definitions --- Config C}
The SQL full-text key and catalog are created by the accepted setup script before the index statement.
\begin{SQLCode}
CREATE FULLTEXT INDEX ON dbo.UC4_Read_C (BODY LANGUAGE 1033)
KEY INDEX PK_UC4_Read_C ON UC4_FT_Read
WITH CHANGE_TRACKING = AUTO, STOPLIST = SYSTEM;
\end{SQLCode}
\begin{MongoCode}
db.uc4_read_C.createIndex(
  {BODY: 'text'},
  {name: 'ix_uc4_body_text', default_language: 'english'}
);
\end{MongoCode}

\ReportBlock{Why It Works}
Looking up a word in an inverted index jumps directly to its posting list in milliseconds instead of scanning millions of characters. \ReportBlock{Trade-Off}
Two costs matter:
\begin{enumerate}
\item \textbf{Semantic Difference vs.\ Substring Matching:} \texttt{LIKE '\%election\%'} matches any character substring (such as \emph{re-election}), whereas linguistic search matches whole words or stems. Thus, baseline and indexed queries do not always return the exact same set of document IDs.
\item \textbf{Async vs.\ Sync Write Maintenance:} Parsing an article adds work. SQL Server performs full-text population in the background with \texttt{CHANGE\_TRACKING = AUTO}, so a new token may become searchable after the write. MongoDB maintains the text index during the write, which can increase write latency.
\end{enumerate}



\UCFourDiagram

\section{UC5: Analytical Aggregation \texorpdfstring{$\rightarrow$}{->} Clustered Columnstore Index (SQL Server)}

\ReportBlock{Problem}
Analytical and reporting queries often scan millions of rows to compute \texttt{GROUP BY} totals across just one or two columns. In a traditional row-oriented table (Rowstore Heap), all 14 columns of each row are stored side by side on 8\,KB pages. To aggregate just \texttt{OBJECT\_TYPE} and \texttt{OBJECT\_ID}, the DBMS must read pages containing columns beyond those needed into memory.

\Needspace{15\baselineskip}
\ReportBlock{Representative Query}
The aggregation groups one million rows by type and reads only \texttt{OBJECT\_TYPE} and \texttt{OBJECT\_ID}. A separate point probe filters \texttt{ID=1} without an ID B-tree on either copy.
\ReportBlock{SQL Server --- aggregation and point probe}
\begin{SQLCode}
SELECT OBJECT_TYPE, COUNT_BIG(*) AS n,
       SUM(COALESCE(OBJECT_ID, 0)) AS s
FROM dbo.UC5_Columnstore GROUP BY OBJECT_TYPE;

SELECT ID, OWNER, OBJECT_NAME, OBJECT_TYPE
FROM dbo.UC5_Columnstore WHERE ID = 1;
\end{SQLCode}

\ReportBlock{Chosen Index}
We compare a 14-column Rowstore Heap (\textbf{Config B}) against a \textbf{Clustered Columnstore Index (CCI)} (\textbf{Config C}) in SQL Server \cite{mssqlcci}. Figure~\ref{fig:diag-uc5} contrasts row and column layouts. A CCI divides rows into rowgroups of up to $1,048,576$ and stores each column in compressed segments. Repeated values in \texttt{OBJECT\_TYPE} make compression promising.

\ReportBlock{Representative index definition --- Config C}
\begin{SQLCode}
CREATE CLUSTERED COLUMNSTORE INDEX CCI_UC5
ON dbo.UC5_Columnstore;
\end{SQLCode}

\ReportBlock{Why It Works}
For the \texttt{GROUP BY} query, SQL Server can read only the two required column segments and process rows in vectorized batches (\textbf{Batch Mode}). Rowstore queries may also use batch mode, so the execution plans must establish which operators used it. \ReportBlock{Trade-Off}
Compressed column segments are not modified in place: small \texttt{INSERT} batches ($<102,400$ rows) can enter an \textbf{\texttt{OPEN} delta rowgroup}, and \texttt{UPDATE}s can mark old rows in a \textbf{delete bitmap} while inserting new versions. This adds work for small writes.

\UCFiveDiagram

\section{UC6: Multi-Tag Membership \texorpdfstring{$\rightarrow$}{->} Multikey vs.\ Normalized Child Index}

\ReportBlock{Problem}
Applications frequently attach multiple categorical tags to each item (for example, tagging a database object with \texttt{owner:SYS}, \texttt{type:TABLE}, and \texttt{temporary:Y}) and search for items that contain either a single tag or a combination of two tags (\texttt{tag1 AND tag2}).

\Needspace{15\baselineskip}
\ReportBlock{Representative Query}
Two single-tag probes use \texttt{type:TABLE} or \texttt{temporary:Y}; two two-tag probes use \texttt{owner:SYS + type:TABLE} or the selective pair below. Both queries return parent IDs.
\ReportBlock{SQL Server --- two-tag membership}
\begin{SQLCode}
SELECT PARENT_ID FROM dbo.UC6_Read_C_Tag
WHERE TAG IN (N'type:TABLE', N'temporary:Y')
GROUP BY PARENT_ID HAVING COUNT(DISTINCT TAG) = 2;
\end{SQLCode}
\ReportBlock{MongoDB --- two-tag membership}
\begin{MongoCode}
db.uc6_read_C.find(
  {tags: {$all: ['type:TABLE', 'temporary:Y']}},
  {_id: 0, PARENT_ID: 1}
);
\end{MongoCode}

\ReportBlock{Chosen Index}
As Figure~\ref{fig:diag-uc6} shows, MongoDB creates one B-tree entry for each distinct tag in a document's array \cite{mongomultikey}. SQL Server stores tags in a normalized child table whose primary key is \texttt{(PARENT\_ID,TAG)}; C adds the reverse key order to seek by tag first.

\ReportBlock{Representative index definitions --- Config C}
The SQL child table retains its \texttt{(PARENT\_ID,TAG)} primary key; C adds the reversed secondary key. MongoDB automatically indexes the distinct array elements as multikey entries.
\begin{SQLCode}
CREATE INDEX IX_UC6_Tag_Parent
ON dbo.UC6_Read_C_Tag (TAG, PARENT_ID);
\end{SQLCode}
\begin{MongoCode}
db.uc6_read_C.createIndex(
  {tags: 1}, {name: 'ix_uc6_tags'}
);
\end{MongoCode}

\ReportBlock{Why It Works}
Both indexes replace full scans with direct B-tree seeks by tag name. \ReportBlock{Trade-Off}
A single MongoDB multikey index entry stores only one tag at a time; for a two-tag \texttt{\$all} query, MongoDB uses the index to seek the more selective tag, fetches those candidate documents from the collection (\texttt{FETCH}), and checks the second tag in memory. Furthermore, if each document has 3 to 5 tags, inserting one document requires inserting 3 to 5 index entries, increasing write overhead.


\UCSixDiagram
\FloatBarrier
